\exercisesection{紧空间上连续函数的代数}

下面所有的 Banach 空间与所有的 Banach 代数都是复的。

\begin{enumerate}
\def\labelenumi{\arabic{enumi})}
\tightlist
\item
  设 X 为拓扑空间，\(\mathcal{C}_{b}(\mathrm{X})\) 为 X 上有界连续复函数所成的星代数，并设 \(\widehat{\mathrm{X}} = \mathrm{X}(\mathcal{C}_{b}(\mathrm{X}))\)。
\end{enumerate}

\begin{enumerate}
\def\labelenumi{\alph{enumi})}
\tightlist
\item
  设 A 为 \(\mathcal{C}_{b}(\mathrm{X})\) 的一个幺 Banach 子代数，设 \(R_{A}\) 为 X 上的如下等价关系：x 与 \(x'\) 等价，如果 \(f(x) = f(x')\) 对每个函数 \(f \in A\) 成立。X 的紧子集 K 的 \(R_{A}\)-饱和集 S 是闭的。（设 \(x \in \overline{S}\)；对一切 \(f \in A\) 及每个 \(\varepsilon > 0\)，设 \(V_{f,\varepsilon}\) 为由那些 \(x' \in X\)（满足 \(|f(x) - f(x')| \leqslant \varepsilon\)）所组成的集合；若 \(f_{1}, \ldots, f_{n} \in A\) 且 \(\varepsilon_{1}, \ldots, \varepsilon_{n} > 0\)，则有
\end{enumerate}

\[
\mathrm{V} _ {f _ {1}, \varepsilon_ {1}} \cap \dots \cap \mathrm{V} _ {f _ {n}, \varepsilon_ {n}} \cap \mathrm{K} \neq \varnothing ,
\]

因此

\[
\bigcap_{\substack{f\in \mathrm{A}\\ \varepsilon >0}}\mathrm{V}_{f,\varepsilon}\cap \mathrm{K}\neq \varnothing
\]

从而 \(x \in S\)。）

\begin{enumerate}
\def\labelenumi{\alph{enumi})}
\setcounter{enumi}{1}
\item
  假设 X 是正规的。设 R 为 X 上的闭等价关系，并设 \(A = A_{R}\) 为由在模 R 的诸类上取常值的函数所组成的 \(\mathcal{C}_{b}(X)\) 的子代数。证明 \(R = R_{A}\)（利用 TG，IX，第 105 页，习题 19，a）。
\item
  假设 X 是紧的。\(R \mapsto A_{R}\) 是从 X 上分离等价关系所成的集合到 \(\mathcal{C}(X)\) 的星子代数所成的集合上的一个双射（设 A 为 \(\mathcal{C}(X)\) 的一个星子代数。证明 \(R_{A}\) 是分离的，并对由 A 过渡到商而得到的 \(X/R_{A}\) 上的连续函数代数应用 Weierstrass–Stone 定理。）
\end{enumerate}

\begin{enumerate}
\def\labelenumi{\arabic{enumi})}
\setcounter{enumi}{1}
\tightlist
\item
  设 X 为紧拓扑空间，R 为 X 上的分离等价关系，B 为 \(\mathcal{C}(\mathrm{X})\) 的包含 \(\mathrm{A}_{\mathrm{R}}\) 的闭子代数（习题 1 的记号）。若 \(f \in \mathcal{C}(\mathrm{X})\) 在每个 R-类 \(c\) 上与 B 的一个元素 \(g_{c}\) 相一致，则 \(f \in \mathrm{B}\)。（设 \(\varepsilon > 0\)。存在 \(\mathrm{V}_{c}\)，即 \(c\) 的一个饱和开邻域，使 \(|f - g_{c}| \leqslant \varepsilon\) 在 \(\mathrm{V}_{c}\) 上成立。设 \(\mathrm{V}_{c_{1}}, \ldots, \mathrm{V}_{c_{n}}\) 为 X 的一个覆盖。设 \((u_{1}, \ldots, u_{n})\) 为从属于 \((\mathrm{V}_{c_{1}}, \ldots, \mathrm{V}_{c_{n}})\) 且由 \(\mathrm{A}_{\mathrm{R}}\) 中的函数组成的单位分解。那么
\end{enumerate}

\[
g = u _ {1} g _ {c _ {1}} + \dots + u _ {n} g _ {c _ {n}} \in \mathrm{B},
\]

且 \(|f - g| \leqslant \varepsilon\)。）

\begin{enumerate}
\def\labelenumi{\arabic{enumi})}
\setcounter{enumi}{2}
\tightlist
\item
  设 X 为紧拓扑空间，且对每个 \(x \in X\)，设 \(A_{x}\) 为交换幺 Banach 代数，用 \(e_{x}\) 表示其单位元。我们假设 \(\|e_{x}\| = 1\)。设 B 为诸 \(A_{x}\)（\(x \in X\)）的乘积 Banach 代数。设 C 为 B 的具有下列性质的 Banach 子代数：
\end{enumerate}

\begin{enumerate}
\def\labelenumi{(\roman{enumi})}
\item
  \(e = (e_x)_{x \in \mathbf{X}} \in \mathbf{C};\)
\item
  若 \(a = (a_x) \in \mathrm{C}\) 且 \(f \in \mathcal{C}(\mathrm{X})\)，则函数 \(fa: x \mapsto f(x)a_x\) 是 \(\mathrm{C}\) 的一个元素；
\item
  对每个 \(a = (a_x) \in \mathrm{C}\)，函数 \(x \mapsto \|a_x\|\) 是上半连续的。
\end{enumerate}

设 \(\widehat{\mathbf{X}}\) 为诸 \(\mathsf{X}(\mathsf{A}_x)\) 的不交并。对每个 \(x_0 \in \mathbf{X}\) 与每个 \(\chi \in \mathsf{X}(\mathsf{A}_{x_0})\)，设 \(\varphi(\chi)\) 为由 \(\varphi(\chi)((a_x)_{x \in \mathbf{X}}) = \chi(a_{x_0})\) 定义的 C 的特征标。 证明 \(\varphi\) 给出从 \(\widehat{\mathbf{X}}\) 到 \(\mathsf{X}(\mathbf{C})\) 上的一个双射。（设 \(\xi \in \mathsf{X}(\mathbf{C})\)。由于 \(\mathcal{C}(\mathbf{X})\) 可等同于 C 的一个闭子代数，\(\xi\) 定义了 \(\mathcal{C}(\mathbf{X})\) 的一个特征标，于是存在 \(x_0 \in \mathbf{X}\)。设 \(a = (a_x) \in \mathbf{C}\) 与 \(b = (b_x) \in \mathbf{C}\)，满足 \(a_{x_0} = b_{x_0}\)；设 \(\varepsilon > 0\)；\(\|a_x - b_x\| < \varepsilon\) 在 \(x_0\) 的某邻域 U 上成立；设 \(f \in \mathcal{C}(\mathbf{X})\) 满足 \(0 \leqslant f \leqslant 1\)，在 \(x_0\) 处等于 1 且在 U 外等于 0；我们有 \(\|fa - fb\| \leqslant \varepsilon\)，从而 \(|\xi(fa - fb)| \leqslant \varepsilon\)，从而 \(|\xi(a) - \xi(b)| \leqslant \varepsilon\)，从而 \(\xi(a) = \xi(b)\)。由此推出存在 \(\chi \in \mathsf{X}(\mathsf{A}_{x_0})\) 使 \(\xi = \varphi(\chi)\)。

\begin{enumerate}
\def\labelenumi{\arabic{enumi})}
\setcounter{enumi}{3}
\tightlist
\item
  设 X 为紧拓扑空间，使得 Banach 代数 \(\mathcal{C}(\mathrm{X})\) 由一个幂等元序列 \((j_n)\) 生成。函数
\end{enumerate}

\[
\omega \mapsto \sum_ {n = 1} ^ {\infty} \frac {1}{3 ^ {n}} (2 j _ {n} (\omega) - 1)
\]

分离 X 的点。它生成 Banach 代数 \(\mathcal{C}(\mathrm{X})\)。

\begin{enumerate}
\def\labelenumi{\arabic{enumi})}
\setcounter{enumi}{4}
\tightlist
\item
  设 A 为从 {[}0,1{]} 到 C 的、在 {[}0,1{]} 上具有直到 n 阶连续导数的函数所成的 Banach 代数（I，第 18 页，例 4）。设 \(\omega \in [0,1] = \mathsf{X}(\mathsf{A})\)。
\end{enumerate}

\begin{enumerate}
\def\labelenumi{\alph{enumi})}
\item
  A 中使 \(V(I) = \{\omega\}\) 成立的最小闭理想 I 是由这样的那些 \(f \in A\) 组成的集合：\(f, f', \ldots, f^{(n)}\) 在 \(\omega\) 处皆为 0。（利用 I，第 94 页，命题 5。）
\item
  若 \(g \in A\)，其典范像在 A/I 中的范数等于
\end{enumerate}

\[
\sum_ {k = 0} ^ {n} \frac {| g ^ {(k)} (\omega) |}{k !}.
\]

\begin{enumerate}
\def\labelenumi{\alph{enumi})}
\setcounter{enumi}{2}
\tightlist
\item
  代数 A/I 同构于 \(\mathbf{C}[\mathrm{X}]/(\mathrm{X}^{n+1})\)；A 中包含 I 的唯一极大理想是诸 \(f \in A\) 中使 \(f(\omega) = 0\) 者所成的集合。
\end{enumerate}

\begin{enumerate}
\def\labelenumi{\arabic{enumi})}
\setcounter{enumi}{5}
\tightlist
\item
  设 \(\Delta\) 为由 \(|\zeta| \leqslant 1\) 给出的 \(\mathbf{C}\) 中的圆盘，A 为在 \(\Delta\) 上连续且在 \(\Delta\) 的内部解析的复值函数所成的 Banach 代数（I，第 20 页，例 9）。
\end{enumerate}

\begin{enumerate}
\def\labelenumi{\alph{enumi})}
\tightlist
\item
  若 \(x \in A\)，则 \(x\) 在 A 中是下列函数的极限：
\end{enumerate}

\[
\zeta \mapsto x _ {n} (\zeta) = x \left(\frac {n \zeta}{n + 1}\right).
\]

\(\Delta\) 的恒等映射是 A 的一个生成元。

\begin{enumerate}
\def\labelenumi{\alph{enumi})}
\setcounter{enumi}{1}
\item
  空间 X(A) 典范地等同于 \(\Delta\)。（利用 I，第 142 页，命题 1 的断言 g）。）
\item
  设 S 为 \(\Delta\) 的一个闭子集，并把它等同于 \(\mathsf{X}(\mathsf{A})\) 的一个闭子集。若 S 有一个非孤立点 \(\zeta\)，满足 \(|\zeta| < 1\)，则 S 对 Jacobson 拓扑的闭包是 \(\Delta\)。
\end{enumerate}

\(d)\) 对 \(x \in \mathrm{A}\)，设 \(x^{*}(\zeta) = x(\overline{\zeta})\)。则 \(x \mapsto x^{*}\) 是 A 的一个等距对合。A 的 Hermite 特征标由 \(\Delta\) 的实元素给出。

\begin{enumerate}
\def\labelenumi{\alph{enumi})}
\setcounter{enumi}{4}
\tightlist
\item
  映射 \(x \mapsto x|U\) 是从 A 到 \(\mathrm{P}(\mathbf{U})\) 上的一个同构。\(\mathcal{C}_{\mathbf{R}}(\mathbf{U})\) 中的每个函数都是 \(\mathrm{P}(\mathbf{U})\) 中函数的实部的一致极限。
\end{enumerate}

\begin{enumerate}
\def\labelenumi{\arabic{enumi})}
\setcounter{enumi}{6}
\tightlist
\item
  设 \(A_{1}\)（相应地，\(A_{2}\)）为习题 5（相应地，习题 6）中所考虑的幺 Banach 代数。设 \(A = A_{1} \times A_{2}\)。证明 A 是半单的，但 A 的 Gelfand 变换的像在 \(\mathcal{C}(\mathsf{X}(\mathsf{A}))\) 中既非闭也非稠密。
\end{enumerate}

¶ 8) 设 X 为紧拓扑空间，B 为 \(\mathcal{C}(\mathrm{X})\) 的一个分离 X 的点的幺 Banach 子代数（带诱导范数）。把 X 等同于 \(\mathrm{X(B)} = \mathrm{X}'\) 的一个闭子集。设 \(\mathrm{B}^*\) 为 B 中可逆元所成的集合。若诸函数 \(\log(|f|)\)（\(f \in \mathrm{B}^*\)）所成的集合在 \(\mathcal{C}_{\mathbf{R}}(\mathrm{X})\) 中稠密，则称 B 为对数模代数（logmodular）。（特别地，当诸函数 \(\mathcal{R}(f)\)（\(f \in \mathrm{B}\)）所成的集合在 \(\mathcal{C}_{\mathbf{R}}(\mathrm{X})\) 中稠密时即是这种情形，因为 \(\log(|e^f|) = \mathcal{R}(f)\)。）以下假设 B 是对数模代数。

\begin{enumerate}
\def\labelenumi{\alph{enumi})}
\tightlist
\item
  对每个 \(\chi \in \mathrm{X}'\)，在 X 上存在唯一的测度 \(\mu_{\chi} \geqslant 0\) 使得
\end{enumerate}

\[
\log (| \chi (f) |) = \int_ {\mathrm{X}} \log (| f (\omega) |) d \mu_ {\chi} (\omega)
\]

对所有 \(f \in \mathrm{B}^*\) 成立，从而 \(\chi(f) = \int_{\mathrm{X}} f(\omega) d\mu_{\chi}(\omega)\) 对所有 \(f \in \mathrm{B}\) 成立。\(b)\) 我们有

\[
\log | \chi (f) | \leqslant \int_ {\mathrm{X}} \log | f (\omega) | d \mu_ {\chi} (\omega)
\]

对所有 \(f \in \mathrm{B}\) 成立，并且条件 \(\chi(f) = \int_{\mathrm{X}} f(\omega) d\mu_{\chi}(\omega)\)（对所有 \(f \in \mathrm{B}\)）唯一地确定 \(\mu_{\chi}\)。

\begin{enumerate}
\def\labelenumi{\alph{enumi})}
\setcounter{enumi}{2}
\tightlist
\item
  设 \(\chi \in X'\)，并设 \(\mu\) 为 \(X\) 上的一个正测度。把 μ 写成 \(\mu = \mu_1 + \mu_2\)，其中 \(\mu_1\) 关于 \(\mu_{\chi}\) 绝对连续，而 \(\mu_2\) 与 \(\mu_{\chi}\) 互相奇异。设 I 为 \(\chi\) 的核。则
\end{enumerate}

\[
\inf _ {f \in \mathrm{I}} \int_ {\mathrm{X}} | 1 - f | ^ {2} d \mu = \inf _ {f \in \mathrm{I}} \int_ {\mathrm{X}} | 1 - f | ^ {2} d \mu_ {1}.
\]

\(d)\) 设 \(\chi\in\mathrm{X}^{\prime}\)，\(h\) 为 \(\mathcal{L}^{1}(\mu_{\chi})\) 中的一个非负函数。设 I 为 \(\chi\) 的核。则

\[
\inf _ {f \in \mathrm{I}} \int_ {\mathrm{X}} | 1 - f | ^ {2} h d \mu_ {\chi} = \exp \Bigl (\int_ {\mathrm{X}} \log (h) d \mu_ {\chi} \Bigr).
\]

\begin{enumerate}
\def\labelenumi{\alph{enumi})}
\setcounter{enumi}{4}
\tightlist
\item
  对 X 上每个正测度 \(\mu\) 以及每个 \(p \in [1, +\infty[\)，设 \(\mathrm{H}^p(\mu)\) 为 B 在 \(\mathrm{L}^p(\mu)\) 中的典范像的闭包。设 \(\chi \in \mathrm{X}'\)，I 为 \(\chi\) 的核。则 \(\mathrm{L}^2(\mu_\chi)\) 是 \(\mathrm{H}^2(\mu_\chi)\) 与在 \(\mathrm{L}^2(\mu_\chi)\) 中所取的 \(\bar{\mathrm{I}}\)（属于 I 的函数的共轭所成的集合）的典范像的闭包的 Hilbert 和。
\end{enumerate}

\(f)\) 空间 \(\mathrm{H}^1 (\mu_\chi)\) 是诸类 \(\widetilde{h}\)（\(h\in \mathcal{L}^1 (\mu_\chi)\)）所成的集合，这些类满足 \(\int_{\mathrm{X}}fh  d\mu_{\chi} = 0\) 对所有 \(f\in \mathrm{I}\) 成立。

\(g)\) 设 \(h\) 为 \(\mathcal{L}^1 (\mu_\chi)\) 中的一个非负函数。则 \(\widetilde{h} = |\widetilde{f}|\)（其中 \(\widetilde{f}\in \mathrm{H}^1 (\mu_\chi)\) 且 \(\int_{\mathrm{X}}f  d\mu_{\chi}\neq 0\)）成立的充分必要条件是 \(\log (h)\in \mathcal{L}^1 (\mu_\chi)\)。

¶ 9) 设 X 为紧拓扑空间，B 为 \(\mathcal{C}(\mathrm{X})\) 的一个幺 Banach 子代数（带诱导范数）。X 的每个元素 \(x\) 都定义 B 的一个特征标 \(\chi_x\)。若 \(x\) 与 \(x'\) 属于 X，则我们记 \(x \sim x'\)，如果 \(\| \chi_x - \chi_{x'} \| \neq 2\)。a) 设 \(x\) 与 \(x'\) 属于 X。若存在 B 中元素的一个序列 \((f_n)\)，使得 \(\| f_n \| \leqslant 1\)、\(|f_n(x)|\) 趋于 1，并且 \(\liminf_{n \to \infty} |f_n(x) - f_n(x')| > 0\)，则 \(x\not\sim x'\)。（可以假设 \(f_n(x)\) 趋于 1。考虑 \(g_n = (f_n - \lambda_n)(1 - \lambda_n f_n)^{-1}\)，其中 \((\lambda_n)\) 是 \([0,1[\) 中趋于 1 的一个数列。我们有 \(\| g_n \| \leqslant 1\)。若序列 \((\lambda_n)\) 选得适当，则 \(g_n(x)\) 趋于 1 且 \(g_n(x')\) 趋于 -1。）

\begin{enumerate}
\def\labelenumi{\alph{enumi})}
\setcounter{enumi}{1}
\item
  设 R 为 B 中元素的实部所成的集合。若 \(x \sim x'\)，则存在 \(c \in ]0,1[\)，使得 \(f(x) \geqslant cf(x')\) 与 \(f(x') \geqslant cf(x)\) 对 R 中每个 \(f \geqslant 0\) 成立。（只需得到第一个条件；若它不成立，则存在 R 中的序列 \((f_n)\)，满足 \(f_n \geqslant 0\)、\(f_n(x') = 1\)，且 \(f_n(x)\) 趋于 0；设 \(g_n \in B\) 满足 \(\mathcal{R}(g_n) = f_n\)。则 \(e^{-g_n} \in \mathrm{B}\)，\(\| e^{-g_n} \| \leqslant 1\)，\(|e^{-g_n(x')}| = 1 / e\)，而 \(|e^{-g_n(x)}|\) 趋于 1，这与 a) 矛盾。）
\item
  在 X 上存在正测度 \(\alpha\) 与 \(\beta\)，使得 \(f(x)-cf(x')=\alpha(f)\) 与 \(f(x')-cf(x)=\beta(f)\) 对 \(f\in R\) 成立。令 \(\mu_{x}=(1-c^{2})^{-1}(c\beta+\alpha)\)，\(\mu_{x'}=(1-c^{2})^{-1}(c\alpha+\beta)\)，则 \(\mu_{x}(f)=f(x)\) 与 \(\mu_{x'}(f)=f(x')\) 对所有 \(f\in B\) 成立，且 \(c\mu_{x}\leqslant\mu_{x'}\)，\(c\mu_{x'}\leqslant\mu_{x}\)。
\item
  仍设 \(x \sim x'\)。证明：若 \(\mu\) 是 X 上的一个正测度，使得 \(\mu(f) = f(x)\) 对所有 \(f \in B\) 成立，则存在 X 上的一个正测度 \(\mu'\)，使得 \(\mu'(f) = f(x')\) 对所有 \(f \in B\) 成立，且 \(c\mu \leqslant \mu'\) 对某个 \(c \in]0,1[\) 成立。（沿用 c) 的记号，令 \(\mu' = \mu_{x'} - c\mu_x + c\mu\)。）
\item
  关系 \(x \sim x'\) 是 X 上的一个等价关系。（利用 d)。）
\end{enumerate}

¶ 10) 设 X 为紧拓扑空间，A 为 \(\mathcal{C}(\mathrm{X})\) 的一个带诱导范数的幺 Banach 子代数。

\begin{enumerate}
\def\labelenumi{\alph{enumi})}
\item
  X 的子集 K 称为（关于 A 的）反对称子集，如果每个在 K 上取实值的函数 \(f \in A\) 都在 K 上取常值。每个反对称子集都包含于某个极大反对称子集。极大反对称子集是闭的。极大反对称子集所成的集合 \(\mathfrak{R}\) 是 X 的一个划分。
\item
  设 \(A^{\perp}\) 为 A 在 X 上复测度所成的 Banach 空间中的正交子空间。设 \(\mu\) 为 \(A^{\perp}\) 的单位球的一个极点。证明 \(\mu\) 的支集是反对称的。
\item
  设 \(f \in \mathcal{C}(\mathrm{X})\)。若 \(f|K \in A|K\) 对所有 \(K \in \mathfrak{R}\) 成立，则 \(f \in A\)。（应用 b) 与 Krein–Milman 定理。）由此重新导出 Weierstrass–Stone 定理（TG，X，第 36 页，定理 3）。
\end{enumerate}

\(d)\) X 的非空子集 E 称为（关于 A 的）峰集（peak），如果存在 \(f\in \mathrm{A}\) 使得 \(\| f\| = 1\) 且 E 是诸 \(x\in \mathrm{X}\) 中使 \(f(x) = 1\) 者所成的集合。于是可以假设 \(|f(x)| < 1\) 对 \(x\notin \mathrm{E}\) 成立（把 \(f\) 换成 \(\frac{1}{2} (1 + f)\)）。峰集的任何非空可数交仍是峰集。若 \(\mathrm{E}_1,\mathrm{E}_2\) 是峰集，则 \(\mathrm{E}_1\cup \mathrm{E}_2\) 是峰集。（设 \(f_{1},f_{2}\in \mathrm{A}\)，使 \(f_{i} = 1\) 在 \(\mathrm{E}_i\) 上成立，\(|f_i| < 1\) 在 \(\mathrm{X - E}_i\) 上成立；设 \(g_{i} = \frac{1}{4} (1 - f_{i})^{1 / 3}\in \mathrm{A}\)；则 \(1 - g_{1}g_{2} = 1\) 在 \(\mathrm{E}_1\cup \mathrm{E}_2\) 上成立，\(|1 - g_1g_2| < 1\) 在 \(\mathrm{X - (E_1\cup E_2)}\) 上成立。）

\begin{enumerate}
\def\labelenumi{\alph{enumi})}
\setcounter{enumi}{4}
\item
  设 E 是峰集的一个交。设 I 为在 E 上为零的 \(f \in A\) 所成的集合。则 \(f \mapsto f|E\) 定义从 A/I 到 A 中函数在 E 上的限制所成的代数 A\textbar E 上的一个等距。
\item
  设 E 为一个峰集。设 \(E' \subset E\) 为关于 A\textbar E 的一个峰集。则 \(E'\) 是关于 A 的一个峰集。
\item
  X 的每个反对称子集都是峰集的一个交。（利用 d) 与 f)。）
\item
  对 X 的每个反对称子集 K，代数 A\textbar K 是 \(\mathcal{C}(\mathrm{K})\) 的一个闭子代数。（利用 \(e\) 与 \(g\)。）
\item
  设 R 为 A 中元素的实部所成的集合。假设 \(X \in \mathfrak{R}\) 且 R 对乘法稳定。则 X 由单独一点组成。（设 \(x_{0} \in X\)。对 \(u \in R\)，存在唯一的 \(f \in A\) 使得 \(u = \mathcal{R}(f)\) 且 \(f(x_{0}) \in \mathbf{R}\)；令 \(\mathrm{N}(u) = \|f\|\)。则 R 对范数 N 是一个实 Banach 空间；闭图像定理证明 R 中的乘法分别连续，从而连续。由此推出 \(S = R + iR\) 上有一个 Banach 代数结构。设 \(p \in R\) 满足 \(p(x) > 0\) 对所有 \(x \in X\) 成立。证明对 S 的每个特征标 \(\chi\) 有 \(\mathcal{R}(\chi(p)) > 0\)，并由此推出 \(\log(p) \in R\)。假设 \(X \neq \{x_{0}\}\)，可以构造一个 \(f \in A\)，使得在 X 上 \(0 < \mathcal{R}(f) \leqslant 1\)，\(f(x_{0}) \in \mathbf{R}\)，且 \(\|f\|\) 可以任意大。根据前文，存在 \(V \in A\) 使 \(|V|^{2} = \mathcal{R}(f)\)。我们有 \(\|V\| \leqslant 1\)，从而 \(\mathrm{N}(\mathcal{R}(V)) \leqslant 2\)，但
\end{enumerate}

\[
\mathrm{N} (\mathcal {R} (\mathrm{V}) ^ {2}) \geqslant \frac {1}{2} \left(\| \mathrm{V} ^ {2} + f \| - | \mathcal {I} (\mathrm{V}) (x _ {0}) ^ {2} |\right)
\]

可以任意大。）

\begin{enumerate}
\def\labelenumi{\alph{enumi})}
\setcounter{enumi}{9}
\tightlist
\item
  若 R 对乘法稳定，则 A = C(X)。（利用 a)、c)、h)、i)。）因此，若 A \(\neq\) C(X)，则存在 \(u \in R\) 使得 \(u^{2} \notin R\)。
\end{enumerate}

\begin{enumerate}
\def\labelenumi{\arabic{enumi})}
\setcounter{enumi}{10}
\tightlist
\item
  设 X 为一度量空间，其距离记作 d。函数 \(f: X \to C\) 称为 Lipschitz 函数，如果存在实数 \(c \geqslant 0\)，使得对 X 中所有 \(x\) 与 \(y\) 有
\end{enumerate}

\[
| f (x) - f (y) | \leqslant c d (x, y)
\]

（参见 FVR，IV，第 7 页，定义 1）。我们用 \(\mathrm{Lip}_b(\mathrm{X})\) 表示 X 上有界 Lipschitz 函数所成的集合。

\begin{enumerate}
\def\labelenumi{\alph{enumi})}
\item
  集合 \(\mathrm{Lip}_b(\mathrm{X})\) 是 \(\mathcal{C}(\mathrm{X})\) 的一个子代数。
\item
  赋以范数
\end{enumerate}

\[
\| f\| = \sup_{x\in \mathrm{X}}|f(x)| + \sup_{\substack{x,y\in \mathrm{X}\\ x\neq y}}\frac{|f(x) - f(y)|}{d(x,y)},
\]

以及对合 \(f \mapsto \overline{f}\) 后，代数 \(\mathrm{Lip}_b(\mathrm{X})\) 是一个交换对合 Banach 代数。它一般不是星代数。

\begin{enumerate}
\def\labelenumi{\alph{enumi})}
\setcounter{enumi}{2}
\item
  对任意交换 Banach 代数 A，我们用 \(\mathsf{X}_{m}(\mathsf{A})\) 表示度量空间 \(\mathsf{X}(\mathsf{A})\)，它赋有由 A 的对偶 \(A'\) 的范数所诱导的距离。度量空间 \(\mathsf{X}_{m}(\mathsf{A})\) 是完备的，且直径 \(\leqslant 2\)。
\item
  设 A 为幺交换 Banach 代数。Gelfand 变换把 A 映到 \(\operatorname{Lip}_{b}(\mathsf{X}_{m}(\mathsf{A}))\) 中。若 A 是半单的，则它诱导从 A 到其像上的一个同胚。
\end{enumerate}

以下我们假设 X 是紧度量空间。记 \(A = \operatorname{Lip}_{b}(X)\)。

\begin{enumerate}
\def\labelenumi{\alph{enumi})}
\setcounter{enumi}{4}
\item
  映射 \(\varphi\) 把 \(x \in X\) 映到特征标 \(f \mapsto f(x)\)，它是从 \(X\) 到 \(X(A)\) 上的一个同胚。
\item
  设 \(d_{1}\) 为由同胚 \(\varphi\) 诱导的 X 上的距离。则 \(d_{1}\) 与 d 等价。
\end{enumerate}

